%% Paper A -- World Patent Information full article.
%% Double-anonymised review requires a separate title page and a de-identified
%% manuscript; the author block is retained here for the production version.
%% cmap must precede fontenc (loaded by cas-dc) so Type-3/pk CM glyphs
%% copy-paste as Unicode rather than dropping the fi/ff ligature.
\RequirePackage{cmap}
\documentclass[a4paper,fleqn]{cas-dc}

\usepackage[numbers]{natbib}
\usepackage{graphicx,booktabs,url}
\input{glyphtounicode}
\pdfgentounicode=1

\graphicspath{{figures/}{../figures/paper_a/}}

\begin{document}
\let\WriteBookmarks\relax
\def\floatpagepagefraction{1}
\def\textpagefraction{.001}

\shorttitle{How much does the counting unit matter?}
\shortauthors{Khan and Israq}

\title[mode=title]{How much does the counting unit matter? Cross-country patent
comparisons and the {WIPO} bulk file}

\author[1]{Samiur Rahman Khan}
\cormark[1]
\ead{samiurk70@gmail.com}
\credit{Conceptualization, Methodology, Software, Data curation, Formal analysis,
Writing -- original draft}
\affiliation[1]{organization={Independent Researcher},
                city={London},
                country={United Kingdom}}

\author[2]{Haseen Israq}
\ead{haseen.israq@g.bracu.ac.bd}
\credit{Validation, Writing -- review and editing}
\affiliation[2]{organization={BRAC University},
                addressline={Pragati Sarani},
                city={Dhaka},
                postcode={1212},
                country={Bangladesh}}

\cortext[1]{Corresponding author}

\begin{abstract}
Open patent bulk files make cross-country comparisons accessible, but their
rows do not necessarily represent inventions. This article measures how
summing the World Intellectual Property Organization (WIPO)
patent-publications-by-technology file across filing offices changes origin
comparisons. For 68 economies with at least 500 publications in 2018--2022,
and after separating the office code that behaves as the PCT international
phase, the office-sum divided by the cell-wise leading office yields a
duplication factor $D$ from 1.05 to 4.48 (median pairwise ratio 1.57). Where
the home office leads, $D=1/H$ arithmetically; where it does not, $D$ stays
far below $1/H$ because foreign filings concentrate in a few large offices.
The fraction of that identity realised, $D\cdot H$, has median 0.90 across
58 measurable-home economies. Partial Spearman of $D$ with European Patent
Convention (EPC) status controlling for $H$ is $+0.413$: the high tail
remains an EPC pattern after home-office share is held. Over seven non-EPC
national-route origins the cell-wise leading office is closer to families
than the office-sum (mean absolute relative error 0.101, or 0.069 excluding
the year-anchor case, versus 0.803). That recommendation is not validated
for the EPC structural case. Holding the family unit fixed, geographic scope
still shifts the China--Singapore comparison from about 500 to about 20.
\end{abstract}

\begin{highlights}
\item Office-summed WIPO publications are not invention counts.
\item Office-sum inflation spans 1.05--4.48 across 68 economies (median pair 1.57).
\item High D remains an EPC pattern after controlling for home-office share.
\item Office \texttt{**} tracks published PCT filings by origin (Spearman $+0.991$).
\item Family scope shifts China--Singapore from about 500x to about 20x.
\end{highlights}

\begin{keywords}
patent statistics \sep patent families \sep WIPO \sep counting unit \sep
PCT \sep cross-country comparison
\end{keywords}

\maketitle

\section{Introduction}\label{sec:intro}

Patent statistics often enter policy reports and cross-country research as if
the word ``patent'' named a single, stable unit. It does not. An application, a
publication, a granted right and a patent family are related records, but they
answer different questions. This distinction has long been central to the
measurement literature because patent data are indicators rather than direct
observations of inventive activity \citep{griliches1990,oecd2009manual}. A
publication count can describe the documentary output of an office. A
first-filing or family count can approximate inventions. A count restricted to
rights sought in several markets can instead indicate an internationally
oriented subset. None is universally preferable; the appropriate unit follows
from the construct being measured.

The problem becomes acute when an analyst works from an aggregate bulk file.
Such files are attractive precisely because they remove many of the barriers
to patent analysis: they are open, compact and already classified by country,
year and technology. Their convenience, however, can hide the joins and
counting rules that would be visible in document-level data. If a table has one
cell for each applicant origin and filing office, summing its office column is
easy. Interpreting that sum as a count of inventions is not. Related
publications protecting one invention can be present at several offices, so
the sum measures publication-to-office assignments. The same procedure can
affect origins differently because filing abroad is not equally common in
every economy.

The WIPO IP Statistics Data Center, and the patent-publications-by-technology
product in particular, is already used for cross-country work. A national
intellectual-property office has produced a regional patent landscape from
this file \citep{ipos2024seasia}. The same Data Center product appears in
applied studies \citep{akkucuk2016patent,hartmann2021georgia,
ijirshar2025innovation,ildiakov2025assessing}. Five uses, and none states
how the office dimension is treated. That is a documented reporting gap, not
an accusation of error. A companion regional analysis uses the family series
for a different substantive purpose \citep{khan2026jseae}; no findings from
that study are required for the measurement presented here.

The established OECD and WIPO counting canon provides the conceptual remedy.
Patent families group related filings; priority-based indicators move closer
to the invention than office publication totals; and carefully selected
international families can improve comparability when domestic institutional
settings differ \citep{oecd2009manual,martinez2010insight,dernis2004triadic,
derassenfosse2013worldwide}. International comparison nevertheless remains
sensitive to attribution, field, year and geographic scope
\citep{hascic2015comparisons}. This literature explains why a unit must be
chosen. It does not tell a user how large the origin-specific inflation is in
the particular open bulk table already on their screen. Comparable source and
classification choices have been shown to matter in applied patent studies,
including work in \emph{World Patent Information} that varies counting rules
in PATSTAT rather than in this bulk file \citep{pasimeni2019assessing}. Here,
the source architecture itself is the object of study.

This article therefore asks a measurement question: how much does summing the
WIPO patent-publications-by-technology bulk file across offices inflate an
origin's count relative to its cell-wise leading office, and does that
inflation behave like a common multiplier across origins? The issue is not
whether every multi-office record can be deduplicated from the aggregate file;
it cannot, because no family identifiers connect the office cells. The issue
is whether a transparent diagnostic can reveal when an office-summed
comparison is especially exposed to filing strategy.

The article makes four contributions. First, it defines what the office-sum
in this WIPO file counts and, equally importantly, what it does not count.
Second, it measures the duplication factor $D$ against the arithmetic identity
$D=1/H$, and it reports that the high tail of $D$ remains associated with
membership of the European Patent Convention (EPC) after home-office share is
controlled. Third, it shows that origin rankings move under
equally defensible counting rules and, where family counts overlap, that the
cell-wise leading office is the closer stand-in for families over the
seven-economy overlap available here. Fourth, holding the family unit
constant, it shows that geographic scope---all families versus
foreign-oriented families---still changes origin comparisons. These
contributions are respectively conceptual, diagnostic, practical and a
matter of geographic scope.

The rest of the article follows the sequence required for a measurement
problem. Section~\ref{sec:units} defines the constructs and the bulk-file
record. Section~\ref{sec:methods} describes the study design, sample and three
quantities, including a worked China--Singapore comparison.
Section~\ref{sec:results} reports the descriptive, residual, rank and
geographic-scope results.
Section~\ref{sec:discussion} discusses their use and limits for patent
information work.

\section{Counting units and the bulk file}\label{sec:units}

\subsection{Families, publications, and office-sums}\label{sec:families}

A patent publication is a document made public by a patent authority at a
particular procedural stage. It is an observable event in one jurisdiction,
not a unique invention. The same underlying inventive subject matter may
generate publications at the applicant's home office, at foreign offices, at
a regional office and through the international phase of the Patent
Cooperation Treaty (PCT). Publication statistics are consequently well suited
to questions about office workload, documentary flows or routes of protection.
They become problematic only when their unit is silently relabelled.

A patent family, by contrast, is a rule-based grouping of related patent
documents. Family definitions vary because different grouping rules serve
different analytical purposes \citep{martinez2010insight}. A simple family may
group documents connected by a common priority; a broader construction can
follow more complex priority links. Triadic families deliberately select
inventions pursued at major offices and are designed for a different
comparability problem \citep{dernis2004triadic}. The important point here is
not that ``family'' is a perfect synonym for invention. It is that a family
definition explicitly groups documents before counting, whereas an office-sum
does not.

WIPO's family indicators used in this article are first-filing based and
attributed by applicant origin \citep{wipo2026families}. Indicator 6a covers
all observed families. Indicator 7 covers families extended to at least one
foreign office. The latter is not a quality-adjusted version of the former and
is not used here to declare one economy more inventive than another. It is a
geographic-scope restriction that selects internationally extended protection.
Such restrictions are common in patent-quality and international-comparison
work because filing in additional markets reflects a selection decision as
well as an opportunity to protect the invention
\citep{squicciarini2013measuring,frietsch2010transnational,
dechezlepretre2017families}.

An office-sum is a third object. For an applicant origin and period, it adds
the publications recorded at each office. If one related set of documents
appears at several offices, each office contributes to the sum. Without
document or family links, the analyst cannot tell which cells relate to the
same invention. The office-sum is therefore neither a simple family nor an
international-family indicator. It is a measure of publication assignments
over an office network.

\subsection{The bulk file as an assignment table}\label{sec:assignment}

The source examined here is WIPO's bulk indicator dataset of patent
publications by publication year, filing office, applicant origin and
technology field \citep{wipo2025bulk}. Each row represents one combination of
those dimensions and supplies a count for that combination. The file is thus
an aggregate assignment table, not a register in which each row is a patent.
The analysed release contains 663,584 rows for 2000--2022 and 50.2 million
publication-to-field assignments.

The assignment-table interpretation matters twice. Across offices, one
invention can contribute publications in several jurisdictions. Across
technology fields, one publication can contribute to several classified
fields. The bulk value is therefore best read as a count attached to a
specific origin--office--year--field cell. Collapsing one dimension does not
automatically create unique documents along another. This is why a large,
cleanly formatted total can still lack the construct implied by the label
``inventions''.

Technology fields follow the Schmoch concordance used for country comparison
\citep{schmoch2008}. The concordance allocates patent classifications to
technology fields, and a document associated with several fields can
contribute to each relevant field. The values in the analysed bulk file are
integers throughout, consistent with whole counting in this particular table.
The aggregate 50.2 million should consequently not be interpreted as the
number of distinct patent documents. For the present study, the technology
dimension is retained consistently in constructing every office quantity, so
the concern is not a technology comparison. It is that the table contains
assignments at more than one level.

This distinction also explains why family counting is not merely a final
arithmetic step on the bulk file. Deduplication requires links. An aggregate
cell reports how many assignments share its dimensions, but it does not report
which assignment in one office corresponds to which assignment elsewhere.
Document-level databases can construct families from priority relations.
This bulk file can only support diagnostics that do not pretend those links
are present.

\subsection{Whole counting and the two \texttt{**} tokens}\label{sec:tokens}

Two identically printed tokens require separate treatment because they occur
in different columns. In the applicant-origin column, \texttt{**} denotes an
unrecorded origin. Such rows cannot be assigned to an economy and are excluded
from origin comparisons. In the filing-office column, \texttt{**} behaves
differently. It is smaller than the named-office sum within a slice and
therefore does not behave as a grand-total field. Its volume is 309,728 in
2020, against WIPO's estimated 275,900 PCT applications filed in 2020
\citep{wipo2021pct}. The two figures differ for two reasons: publications lag
applications by about eighteen months, and whole-counted field assignments
exceed document counts. The distribution of office \texttt{**} is higher for
origins that make heavier use of international filing.

WIPO's Data Center documentation states that international-phase PCT data are
available under the PCT route, whereas subsequent national-phase activity is
reported under patent or utility-model indicators \citep{wipoabout}. In
combination with the scale and origin pattern of office \texttt{**}, this
supports treating that office token as the PCT international phase.
Section~\ref{sec:pct-result} reports a discriminating origin-level rank check
against published PCT applications by origin. The mapping is nevertheless
inferential: the documentation available for the Data Center does not
explicitly define the \texttt{**} value in this bulk file's office column.
This article therefore says that office \texttt{**} \emph{behaves as} the PCT
layer. It does not claim that the token's meaning has been directly confirmed
by a file-specific data dictionary.

Separating the token is necessary for a coherent office-sum. If it were added
to all named offices, an international-phase publication would be treated as
another office observation beside later national or regional publications.
That may be appropriate when studying the use of the PCT route, but not when
the stated outcome is publications at named patent offices. The origin token
and office token are therefore removed or separated for different reasons:
one lacks geographic attribution; the other represents a distinct filing
layer under the maintained interpretation.

Figure~\ref{fig:schematic} states the counting problem as two illustrative
filing routes. It is not an average family. The internationally filed route
returns six, five, one, one and one under the five rules shown; the
domestic-only route returns one, one, one, one and zero. Inflation is
therefore a property of the route, not a common multiplier, and indicator 7
returns zero for a domestic invention rather than a smaller positive number.

\begin{figure*}[t]
\centering
\includegraphics[width=0.95\textwidth]{figA0_counting_schematic}
\caption{Two illustrative filing routes, not an average family size. Panel A
follows a priority filing through a PCT international-phase publication,
three foreign offices and a regional (EPO) publication. Panel B is a
home-office publication only. The five counts under each panel are
office-sum including \texttt{**}, office-sum excluding \texttt{**}, cell-wise
leading office, all families (indicator 6a) and foreign-oriented families
(indicator 7): $6/5/1/1/1$ versus $1/1/1/1/0$. The time axis marks priority,
publication (eighteen months) and national-phase entry (thirty months).
Source: authors' schematic.}
\label{fig:schematic}
\end{figure*}

\section{Materials and methods}\label{sec:methods}

\subsection{Study design and unit of analysis}\label{sec:design}

The study is a descriptive measurement analysis of applicant-origin filing
profiles. Its unit of comparison is an economy over the 2018--2022 window.
The design does not estimate causal effects and does not rank national
innovation systems. Instead, it asks whether the office distribution attached
to one applicant origin makes an office-summed publication count more or less
concentrated than the distribution attached to another origin.

The primary calculation begins with origin--year--technology-field cells.
Named offices are compared within each cell, and the resulting counts are
aggregated over the study window. Retaining a common cell structure prevents a
country's technology composition or year mix from being silently collapsed at
a different stage from another country's. Office \texttt{**} is separated
before the named-office duplication proxy is calculated. Rows with origin
\texttt{**} are excluded because they cannot enter an origin-specific result.

Three descriptive quantities are then placed in sequence. The home-office
share describes how much of an origin's named-office activity is observed at
its own office. The PCT-layer share describes the relative size of the
separated \texttt{**} office route. The duplication factor describes how
dispersed the named-office count is relative to the cell-wise leading office.
Together they distinguish domestic concentration, use of the inferred PCT
route and multi-office dispersion. None directly observes family links.

\subsection{Data sources and pin}\label{sec:sources}

Publication counts come from the WIPO patent-publications-by-technology bulk
file, retrieved 30 July 2026 \citep{wipo2025bulk}. The cached input is pinned
in the project manifest by SHA-256 so that the analysed release can be
distinguished from later updates. The retrieval, transformation and figure
scripts preserve the source file and generate the derived origin-level table.
This pin is important for a live statistical service: a retrieval date alone
identifies when access occurred, whereas the checksum identifies the exact
bytes analysed.

Family counts come from a documented manual export of WIPO Data Center
indicators 6a and 7, total count by applicant's origin
\citep{wipo2026families}. The Data Center recorded a database update in May
2026, and the export was retrieved 2 September 2026. The export route is
manual because the family series is exposed through the Data Center interface
rather than the publication bulk download. As a provenance check, the saved
series reproduces, within source rounding, the China and Japan values printed
for the same indicators in \emph{World Intellectual Property Indicators 2024}
\citep{wipo2024wipi}.

The two data products are not merged record by record. The publication file is
dated by publication year and the family series by earliest filing year.
Their annual observations therefore have different temporal anchors, and one
series is not assumed to nest inside the other. WIPO's Data Center
documentation states that both the technology indicators and the family
indicators are compiled from PATSTAT \citep{wipoabout}. The two products
therefore share a lineage, so the non-nesting of annual observations is a
dating-rule artefact rather than a difference of source. The family export is
used first for a window-level consistency check and then for a separate
geographic-scope illustration. These uses avoid presenting a same-year
difference between the products as though it were a count of duplicates.

\subsection{Sample construction}\label{sec:sample}

Sample construction begins with all 663,584 bulk rows. The analysis retains
positive technology-field codes, excludes unrecorded applicant origins and
forms origin-specific totals over 2018--2022. Economies enter the comparison
when they have at least 500 publications in that window. The resulting sample
contains 68 economies. The threshold screens volume, not attribution
meaningfulness. Repeating the range calculation with a threshold of 200 or
1,000 leaves the observed origin--year--field duplication-factor range at
1.05--4.48.

The home-office share requires a correspondence between applicant origin and
an origin-resolved office row. This is measurable for 58 economies. For 10
economies, the bulk file has no origin-resolved row for the corresponding home
office. That absence is not coded as a zero. A missing home-office block can
reflect source coverage rather than no resident filing, so plotting it at zero
would attach a behavioural interpretation that the data do not support. These
economies remain in duplication and PCT calculations based on the observed
offices, but they are excluded from the domestic-share distribution and from
the horizontal axis of Fig.~\ref{fig:dup}.

$D$ is a property of the cell, not only of the file. The direction of
coarsening is arithmetic: the sum of within-cell maxima is at least the
maximum of the summed cells, so coarser cells cannot lower $D$. The magnitude
is the finding. At origin--year--field granularity the range is 1.05 / 2.27 /
4.48. Collapsing the field dimension raises it to 1.05 / 2.51 / 5.41 at
origin--year cells. Collapsing year while retaining field isolates the
temporal-spread component: origin--field $D$ is 1.05 / 2.39 / 4.90, so the
maximum rises from 4.48 to 4.90 and the median from 2.27 to 2.39. Origin-only
$D$ is 1.05 / 2.54 / 5.77. That origin-only quantity is the office-sum
divided by the largest single named office---the definition previously
labelled and then corrected away from. Coarser cells therefore raise $D$; they
do not compress the 4.3-fold headline. This article reports the
origin--year--field $D$ as the main measure and the coarser range as a check.
A missing home-office block can still raise the calculated proxy if the
unobserved office would have been the cell-wise leader.

\subsection{Three quantities and a worked example}\label{sec:quantities}

Let $p_{ioyt}$ denote the publication count for applicant origin $i$, named
office $o$, publication year $y$ and technology field $t$. Let $h(i)$ be the
home office corresponding to origin $i$. The home-office share is
\[
 H_i =
 \frac{\sum_{y,t}p_{i,h(i),y,t}}
      {\sum_{o,y,t}p_{ioyt}},
\]
where the denominator contains named offices only. The measure is undefined
when the source lacks an origin-resolved block for $h(i)$. It is not replaced
by zero.

Let $q_{iyt}$ denote the count under office \texttt{**}. The PCT-layer share,
under the inferential identification in Section~\ref{sec:tokens}, is
\[
 Q_i =
 \frac{\sum_{y,t}q_{iyt}}
      {\sum_{o,y,t}p_{ioyt}+\sum_{y,t}q_{iyt}}.
\]
This denominator is deliberately different from that of the home-office
share: $Q_i$ asks how large the separate layer is when it is placed beside the
named-office total.

Finally, define the duplication factor as
\[
 D_i =
 \frac{\sum_{y,t}\sum_o p_{ioyt}}
      {\sum_{y,t}\max_o(p_{ioyt})}.
\]
The numerator is the named-office sum. The denominator adds the cell-wise
leading-office count within each origin--year--technology-field cell. That
denominator is the sum of within-cell maxima. It is therefore less than or
equal to the origin-level maximum office, $\max_o(\sum_{y,t}p_{ioyt})$; the
two denominators differ for 61 of 68 economies (median gap 0.224, maximum
2.294 at Monaco). The origin-only grain in Section~\ref{sec:sample} is exactly
office-sum divided by that origin-level maximum, which is why the denominators
differ. Computed $D$ is the lower of the two ratios. The factor is bounded
below by unity and rises as meaningful volume is spread across additional
offices. A regional office, including the European Patent Office, enters as
one observed office; no equivalent-count multiplier is applied.

The duplication factor confounds office dispersion with the temporal spread
of publication, because related publications that land in different years
occupy separate cells and therefore add separate cell-wise maxima. The
origin--field grain, collapsing year while retaining field, isolates that
timing component: the maximum rises from 4.48 to 4.90 and the median from
2.27 to 2.39.

When the home office is the cell-wise leader, $D=1/H$ exactly, up to cells
where another office leads. China, Korea and Japan sit on that identity
($H=95.4\%$, $68.1\%$ and $55.4\%$; $D=1.05$, $1.47$ and $1.80$). Singapore's
home office is not the leader---the United States office is---so the identity
does not bind: $H=5.0\%$, $1/H=20.08$, $D=2.60$ and residual $17.48$.

China provides the domestic-concentration case. In 2018--2022, 95.4\% of its
named-office publications are at its home office, and its duplication factor
is 1.05. The worked reading is straightforward: within the bulk file, the
named-office sum is only modestly larger than the contribution of the
cell-wise leading office. This does not prove that the leading-office count is
a family count. It shows that adding other observed offices makes little
difference to this particular office-based total.

Singapore provides the contrasting filing profile. Its home-office share is
5.0\%, so its own office cannot stand in for most publications attributed
to Singaporean applicants in the file. Its activity is distributed across
foreign offices, and the office-sum accumulates those observations. The
contrast with China is therefore not an arithmetic nuisance that a common
percentage adjustment can remove. The two origins place their publications
through different office networks.

The sample maximum further clarifies the denominator. Switzerland has a
duplication factor of 4.48, a home-office share of 3.1\%, and its
origin-level leading office is the EPO rather than its home office. The
factor does not say that each Swiss invention has 4.48 publications, nor does
it place an exact bound on distinct inventions. It says that, under the
cell-wise construction, the named-office sum is much larger relative to the
leading-office baseline than it is for China. The measure is consequently
labelled a duplication \emph{proxy}.

\subsection{Check against WIPO family indicators}\label{sec:check}

A useful proxy should at least order filing profiles similarly to a measure
that uses families, even though the levels need not coincide. The check
compares $D_i$ with publications divided by WIPO indicator 6a families over
the common 2018--2022 window. Seven economies overlap: China, Japan, Korea,
Malaysia, the Philippines, Singapore and Thailand---all non-EPC origins that
file through national rather than regional-first routes. Spearman correlation
is used because the diagnostic claim concerns ordering rather than a linear
conversion.

The two ratios are not independent measures. $D$ is the office-sum divided by
the cell-wise leading office; publications per family uses the same office-sum
as its numerator. Rank agreement is therefore expected to be high if the
mechanism exists at all. The practical comparison is whether the
leading-office count or the office-sum lies closer to families. Publications
and families use different dating rules, and $D$ cannot fall below unity
whereas the observed publications-to-families ratio can. The exercise is a
consistency check on behaviour, not an estimated deduplication formula.
Family exports for four EPC origins were not available in the pinned input.
Section~\ref{sec:proxy-result} reports the coefficients and the error
comparison.

\subsection{Geographic-scope comparison: indicators 6a and 7}\label{sec:scope-method}

The final analysis holds the family unit constant and changes geographic
scope. WIPO indicator 6a counts all families by applicant origin and earliest
filing year. Indicator 7 counts the subset extended to at least one foreign
office. China and Singapore are used because their domestic shares in the
bulk file already establish sharply contrasting filing profiles. Their
family series from 2000--2022 show whether the substantive comparison also
depends on including or excluding domestic-only families.

The all-family and foreign-oriented ratios are not alternatives from which the
article selects a preferred headline. They answer different questions. The
first describes the full first-filing-based family portfolio observed for an
origin. The second describes the internationally extended subset. Placing
them together demonstrates that correcting the unit does not eliminate the
need to define scope.

\section{Results}\label{sec:results}

\subsection{Sample and range table}\label{sec:range}

Table~\ref{tab:range} summarises the three bulk-file quantities. Across the 58
economies whose home office is measurable, the home-office share ranges from
0.3\% (Saudi Arabia) to 95.4\% (China), with a median of 37.8\% (p25 9.2,
p75 64.8). Across all 68 economies, the inferred PCT-layer share ranges from
3.2\% (Macao SAR) to 54.4\% (Iran) and has a median of 15.1\% (p25 12.7,
p75 18.6). Both PCT endpoints are origins whose home office is unmeasurable.
The duplication factor ranges from 1.05 to 4.48. Its all-economy median is
2.27 (p25 1.55, p75 3.34); among the 58 economies with a measurable home
office, it is 2.14 (p25 1.48, p75 3.45).

\begin{table*}[t]
\caption{Filing-profile measures in the WIPO bulk file, 2018--2022.
The sample contains 68 economies with at least 500 publications; the
home-office row contains the 58 economies for which the own office is
origin-resolved. The fourth row reports $D$ on that measurable-home base.
The fifth row is the fraction of the identity $D=1/H$ actually realised.
Per-economy values are in \texttt{counting\_generality.csv}
in the replication archive. Source: authors' calculations from the WIPO bulk
patent-indicators file.}
\label{tab:range}
\centering
\small
\begin{tabular}{@{}lcccccc@{}}
\toprule
Measure & $n$ & Min. & p25 & Median & p75 & Max. \\
\midrule
Home-office share (\%) & 58 & 0.3 & 9.2 & 37.8 & 64.8 & 95.4 \\
PCT-layer share (\%) & 68 & 3.2 & 12.7 & 15.1 & 18.6 & 54.4 \\
Duplication factor & 68 & 1.05 & 1.55 & 2.27 & 3.34 & 4.48 \\
Duplication factor (measurable home) & 58 & 1.05 & 1.48 & 2.14 & 3.45 & 4.48 \\
$D\cdot H$ (measurable home) & 58 & 0.005 & 0.32 & 0.90 & 0.99 & 1.00 \\
\bottomrule
\end{tabular}
\end{table*}

The duplication proxy does not apply a common uplift across origins. The more
robust summary is the median pairwise ratio of 1.57, the median of the 2,278
pairwise ratios $D_i/D_j$ taken larger over smaller. The 4.3-fold figure is
the ratio of the sample ends, Switzerland versus China. These are differences
in $D_i$, not estimates of the number of copies attached to an invention.
They establish the empirical range over which the same office-summing
operation behaves differently in this file. The pooled median of 2.27 is a
composition artefact: 35 of the 68 are EPC contracting-state origins. Group
medians on the 58 economies with a measurable home office are the more stable
summary.

\subsection{Residual from the identity $D=1/H$}\label{sec:dup-result}

Figure~\ref{fig:dup} places the duplication factor against the home-office
share for the 58 measurable economies, with the curve $D=1/H$ overlaid.
Nineteen of those 58 lie within 0.05 of the identity. Thirty-five of 58 have
a home office that is the origin-level leading named office. All 19 are among
those 35. The remaining 16 lead at origin level but not in every year-field
cell, so cell-wise $D$ sits below $1/H$. China occupies the low-duplication,
domestic-heavy end at 95.4\% and 1.05, where the identity binds. The
information in the scatter is not a downward-sloping cloud. It is the gap
below $1/H$. At low home-office share the identity predicts $D$ far above the
observed ceiling.

The residual $1/H-D$ inherits that explosion: its maximum is set by how small
the smallest $H$ happens to be. Across the 58 it runs from 0.00 to 369.68.
The scale-free counterpart is $D\cdot H$, the fraction of the identity
actually realised, bounded in $(0,1]$. It has median 0.90 (p25 0.32, p75
0.99); 23 of 58 realise at least 0.95 of it. Economies on the identity
realise essentially all of it (China, Korea and Japan at 1.00; Malaysia at
0.96). The low-$H$ group does not: Singapore 0.13, Switzerland 0.14, Saudi
Arabia 0.005. Economies that file away from home do not disperse across all
offices. They concentrate in a handful of large ones, so office-sum inflation
is bounded far below what home-office share alone would imply.

Saudi Arabia has $H=0.3\%$, $1/H=371.41$, $D=1.73$ and residual $369.68$, with
32 origin-resolved home-office publications. That count places it at the
boundary of the measurable/unmeasurable distinction rather than clearly inside
it: a partially resolved home-office block also produces a small non-zero
count, and 32 publications over five years cannot carry the paper's largest
single number as a settled filing pattern. Among economies with at least
1,000 origin-resolved home-office publications, the largest residual is
Israel: $H=2.9\%$, $1/H=34.95$, $D=2.55$ and residual $32.40$, on 1,450 home
publications. Singapore remains the interpretable mid-case: $1/H=20.08$,
$D=2.60$ and residual $17.48$.

\begin{figure*}[t]
\centering
\includegraphics[width=0.88\textwidth]{figA1_duplication_vs_domestic}
\caption{Duplication factor versus the applicant-origin home-office share,
2018--2022, for the 58 economies with a measurable home office. The dashed
curve is the identity $D=1/H$, which is arithmetic when the home office is
the cell-wise leader. Marker fill denotes whether the home office is the
origin-level leading named office; colour denotes EPC versus non-EPC origin.
At low $H$, Singapore has $1/H=20.08$ and $D=2.60$; Saudi Arabia has
$1/H=371.41$ and $D=1.73$; the observed ceiling is 4.48. Ten economies remain
in the range statistics but are omitted here because the home office is not
origin-resolved. Source: authors' calculations from the WIPO bulk
patent-indicators file.}
\label{fig:dup}
\end{figure*}

The high tail of $D$ remains an EPC pattern after home-office share is held.
Partial Spearman of $D$ with EPC status controlling for $H$ is $+0.413$
($p=0.0016$; 95\% CI $[0.171, 0.608]$) across the 58. Table~\ref{tab:epcq}
shows where that association sits: in the lowest $H$ quartile, median $D$ is
2.57 among six non-EPC origins and 3.77 among nine EPC origins. The gap
narrows in the upper half of $H$. The second quartile's non-EPC cell has only
two origins and is not a basis for a separate claim. At the lowest $H$ the
values still interleave---New Zealand 3.87 beside Sweden 3.88---and none of
those four matched-$H$ origins is led by the EPO (all four are led by the
United States office). EP-as-leader is therefore not required for high $D$.
The matched comparison does not replace the partial correlation; it shows
that the residual EPC association is a shift in the low-$H$ distribution, not
a requirement that the regional office lead.

\begin{table}[t]
\caption{Median duplication factor by home-office-share quartile and EPC
status, 2018--2022, measurable-home base ($n=58$). Quartiles are of $H$.
The second-quartile non-EPC cell has $n=2$. Source: authors' calculations
from the WIPO bulk patent-indicators file.}
\label{tab:epcq}
\centering
\small
\setlength{\tabcolsep}{4pt}
\begin{tabular}{@{}lcccc@{}}
\toprule
Quartile & $n$ & non-EPC $D$ & $n$ & EPC $D$ \\
\midrule
Q1 (low $H$) & 6 & 2.57 & 9 & 3.77 \\
Q2 & 2 & 2.84 & 12 & 3.67 \\
Q3 & 7 & 1.80 & 7 & 2.01 \\
Q4 (high $H$) & 11 & 1.23 & 4 & 1.33 \\
\bottomrule
\end{tabular}
\end{table}

\begin{table}[t]
\caption{Duplication factor by EPC status, 2018--2022. The measurable-home
base of 58 is the primary comparison; the all-68 rows include ten origins
whose home office is not origin-resolved. Source: authors' calculations from
the WIPO bulk patent-indicators file.}
\label{tab:epc}
\centering
\small
\begin{tabular}{@{}llcccc@{}}
\toprule
Base & Group & $n$ & Min. & Median & Max. \\
\midrule
58 & EPC & 32 & 1.15 & 2.82 & 4.48 \\
58 & non-EPC & 26 & 1.05 & 1.63 & 3.87 \\
68 & EPC & 35 & 1.15 & 2.75 & 4.48 \\
68 & non-EPC & 33 & 1.05 & 1.73 & 4.10 \\
\bottomrule
\end{tabular}
\end{table}

Table~\ref{tab:epc} records the group ranges. On the measurable-home base
the non-EPC maximum is New Zealand. The five highest values on that base are all
EPC: Switzerland, Luxembourg, Norway, Belgium and the Netherlands. Ten of the
68 origins have the EPO as origin-level leading office; all ten are EPC
origins. The all-68 non-EPC maximum is Antigua and Barbuda, whose home office
is unmeasurable.

\subsection{The PCT layer}\label{sec:pct-result}

The separated \texttt{**} office layer also varies by applicant origin.
Its sample range is 3.2\% (Macao SAR) to 54.4\% (Iran), with a median of
15.1\%. Iran routing over half of its publications through office
\texttt{**} is the sample maximum; because its home office is unmeasurable,
that endpoint is taken up with Antigua and Monaco in
Section~\ref{sec:limitations} rather than treated here as a confirmed filing
behaviour. In the China--Singapore comparison, the layer accounts for 3.9\%
and 20.2\%, respectively, over 2018--2022. Adding this layer to named-office
publications would therefore change the two origin totals by different
proportions.

If office \texttt{**} is the PCT international phase, the PCT-layer share
$Q$ should rise as the home-office share $H$ falls. Over the 58 measurable
economies, Spearman correlation of $Q$ with $H$ is $-0.559$ ($p=0.000005$;
95\% CI $[-0.714, -0.352]$). The relationship is descriptive, not
confirmation: $H$ and $Q$ share the named-office sum in different roles in
their denominators, and the scatter is fat. It also does not discriminate
against the rival reading. If \texttt{**} were an unresolved-office bucket,
unresolved-office publications would still come disproportionately from
origins whose filings are dispersed, which is the same low-$H$ group. Both
readings imply a negative $Q$--$H$ correlation.

A rank check against published PCT applications by origin does discriminate.
If office \texttt{**} is the international phase, origin-level \texttt{**}
volumes should track WIPO's published PCT filings; an unknown-office bucket
would have no reason to. Spearman correlation of 2020 office-\texttt{**}
publications against published 2020 PCT applications by origin is $+0.991$
over the twenty largest named origins ($p<10^{-15}$; 95\% CI $[0.977,
0.997]$), using the same 2020 office-\texttt{**} construction as the 309,728
aggregate \citep{wipo2021pct}. Sweden matches exactly: 4,356 publications
under \texttt{**} and 4,356 published PCT applications. The $Q$--$H$
correlation is consistent with the PCT identification without ruling out the
alternative; this origin-level check is the test that would have failed if
\texttt{**} were merely unresolved offices.

Table~\ref{tab:pct} reports the PCT-layer share for selected origins in the
68, after that correlation. The values use the same $\texttt{tec\_id}>0$,
2018--2022 and 68-economy filters as the rest of the paper; they are not the
all-year audit table. Russia (4.5\%) and China (3.9\%) sit below Korea
(8.8\%) and India (7.5\%), which in turn sit below Japan (14.3\%), Germany
(13.7\%) and the United States (13.3\%); the Netherlands (16.6\%) and
Sweden (18.8\%) are higher still.

\begin{table}[t]
\caption{PCT-layer share of publications for selected origins in the 68,
2018--2022, $\texttt{tec\_id}>0$, under the same 68-economy filters as the
range table; not the all-year audit table. Office \texttt{**} is treated as
the PCT international phase on behavioural grounds and on the origin-level
rank check in Section~\ref{sec:pct-result}; the mapping is still without a
file-specific dictionary. Source: authors' calculations from the WIPO bulk
patent-indicators file.}
\label{tab:pct}
\centering
\small
\begin{tabular}{@{}lc@{}}
\toprule
Origin & PCT-layer share (\%) \\
\midrule
Russia & 4.5 \\
China & 3.9 \\
Korea & 8.8 \\
India & 7.5 \\
Japan & 14.3 \\
Germany & 13.7 \\
United States & 13.3 \\
Netherlands & 16.6 \\
Sweden & 18.8 \\
\bottomrule
\end{tabular}
\end{table}

This result is distinct from Fig.~\ref{fig:dup}. The duplication factor
describes dispersion across named offices after the token has been separated.
The PCT-layer share describes the size of the removed route relative to the
combined total. An analyst interested in PCT use may retain and report it as
an outcome. An analyst describing national or regional office publications
should separate it and state that treatment. Neither choice converts
publications into families.

\subsection{Rank instability}\label{sec:ranks}

Figure~\ref{fig:ranks} ranks the 25 largest origins by named-office
publication volume under three rules, each legitimate for a different
question: the office-sum including \texttt{**}, the office-sum excluding
\texttt{**}, and the cell-wise leading office. None is the correct ranking.
The figure is restricted to the top 25 by named-office volume because the
full 68 is unreadable, not in order to hide a larger movement outside that
window. Among those 25, Russia ranks 13th on the office-sum including
\texttt{**} and 6th on the cell-wise leading office, because its PCT-layer
share is 4.5\% and its domestic home-office share is 91.9\%, so removing
\texttt{**} and moving to the leading office penalises internationally
dispersed origins more than Russia. T\"urkiye ranks 23rd versus 15th; Brazil
24th versus 16th; Poland 26th versus 18th. The largest movement among the 68
is Antigua and Barbuda, 50th versus 65th, a first-named-applicant domicile
case taken up in Section~\ref{sec:limitations}. The figure is a consequence
of the range in Table~\ref{tab:range}: if $D$ is not a common multiplier,
origin order is not invariant to the counting rule.

\begin{figure*}[t]
\centering
\includegraphics[width=0.88\textwidth]{figA5_rank_instability}
\caption{Rank of the 25 largest origins by named-office publication volume
under three counting rules, 2018--2022. The columns are the office-sum
including \texttt{**}, the office-sum excluding \texttt{**}, and the
cell-wise leading office. All three are legitimate answers to different
questions; none is the correct ranking. The colour bar reaches 8. Among this
top 25, Russia ranks 13th (sum including \texttt{**}) versus 6th (cell-wise
leading office); T\"urkiye 23rd versus 15th; Brazil 24th versus 16th; Poland
26th versus 18th. The figure is restricted to the top 25 because the full 68
is unreadable. Antigua and Barbuda, 50th versus 65th, is the largest
movement among the 68. Source: authors' calculations from the WIPO bulk
patent-indicators file.}
\label{fig:ranks}
\end{figure*}

\subsection{Proxy versus families}\label{sec:proxy-result}

Figure~\ref{fig:proxy} places WIPO all-family counts on a logarithmic
horizontal axis against the cell-wise leading office (circles) and the
office-sum (squares) for the seven overlapping economies. The dashed line is
1:1. The leading-office series tracks that line; the office-sum series lies
above it, except where dating rules pull a point toward or below equality.
The Philippines office-sum can sit near or below the 1:1 line because
publication year and earliest filing year do not nest; that is a feature of
the dating rules in Section~\ref{sec:sources}, not an anomaly in the
Philippines series.

\begin{figure}[t]
\centering
\includegraphics[width=\columnwidth]{figA2_proxy_validation}
\caption{WIPO patent families (indicator 6a) on a logarithmic horizontal axis
against the cell-wise leading office (circles) and the office-sum (squares),
2018--2022, for the seven economies available in both sources. The dashed
line is 1:1. A faint vertical connector joins each pair; that gap is that
economy's $D$. The Philippines office-sum can sit near or below the 1:1 line
because publication year and earliest filing year do not nest. Source:
authors' calculations from the WIPO bulk patent-indicators file and WIPO
indicator 6a.}
\label{fig:proxy}
\end{figure}

Spearman $+1.000$ between the cell-wise leading office and families on $n=7$
is a size ranking, not an independent test. The office-sum against families
returns Spearman $+0.964$ ($p=0.0005$; 95\% CI $[0.771, 0.995]$), the same
coefficient as $D$ against publications/families, because those ratios share
the office-sum numerator. Rank agreement barely discriminates the two
candidates. What separates them is the mean absolute relative error: $0.101$
for the leading office versus $0.803$ for the office-sum, or $0.069$ versus
$0.803$ once the Philippines is excluded. The year-anchor offset---publication
year versus earliest filing year---applies to all seven economies. It is
visible in the Philippines because that origin's high home-office share puts
its ratio closest to unity to begin with (leading-office/families $0.70$).
Singapore has $|\mathrm{lead}/\mathrm{fam}-1|=0.22$, the largest remaining
error. Spearman
correlation of (leading-office/families) with $D$ is $-0.214$ ($p=0.6445$):
there is no evidence at $n=7$ that $D$ predicts the error a user would incur,
and $D$ is not a validated warning light for that proxy error. Seven points
support a scoped recommendation for this overlap, all non-EPC origins filing
through national rather than regional-first routes. They do not support a
general calibration equation, they do not validate the recommendation for the
EPC structural case, and they do not support dividing an office-sum by $D_i$
and calling the result families.

\subsection{China--Singapore on all families versus foreign-oriented families}
\label{sec:scope-result}

Figure~\ref{fig:scope} reports the China-to-Singapore ratio while holding the
family unit constant. China's all-family series rises from 23,079 in 2000 to
1,517,232 in 2022; Singapore's rises from 716 to 3,036. On all families, the
2022 ratio is about five hundred. When the scope is restricted to families
extended abroad, the ratio is about twenty. The vertical axis is logarithmic
so that the foreign-oriented series remains readable. The 2019 dip in the
all-family series is in the data (ratio 416). Values for 2022 are
provisional because later foreign members may not yet be observed.

\begin{figure*}[t]
\centering
\includegraphics[width=0.88\textwidth]{figA4_china_singapore_ratio}
\caption{China-to-Singapore patent-family ratio under two geographic scopes,
by earliest filing year, 2000--2022, on a logarithmic vertical axis. The
solid series uses all families (WIPO indicator 6a); the dashed series uses
only families extended to at least one foreign office (indicator 7). The
2019 dip in the all-family series is in the data (ratio 416). Values for
2022 are provisional because later foreign members may not yet be observed.
Source: authors' calculations from WIPO IP Statistics Data Center indicators
6a and 7.}
\label{fig:scope}
\end{figure*}

The difference corresponds to the composition of the two portfolios. In
2022, 3.6\% of Chinese families are observed as extended abroad, compared with
91.9\% of Singaporean families. Because recent first filings may acquire
foreign members after their earliest filing year, the 2022
foreign-oriented result is provisional. The more mature 2018 anchor gives
the same qualitative contrast: about 470 on all families and about 15 on
foreign-oriented families.

\section{Discussion}\label{sec:discussion}

\subsection{What an office-sum measures}\label{sec:meaning}

Three findings do not depend on the scoped checks that follow. First, the
office-sum is not an invention count. Relative to the cell-wise leading
office, inflation ranges from 1.05 to 4.48 across 68 economies, with a median
pairwise ratio of 1.57; it is not a common multiplier. Second, where the home
office leads, $D$ is arithmetically $1/H$, so most observed variation in $D$
is a restatement of home-office share. Where it does not lead, $D$ stays far
below $1/H$ because foreign filings concentrate in a few large offices, and
$D\cdot H$ is the fraction of that identity the data deliver. Third, origin
rankings are not invariant to defensible counting rules, demonstrated on the
top 25. Everything else in this section is scoped.

The office-sum is not meaningless; it is misnamed when presented as inventions.
Properly labelled, it describes the volume of publication assignments
generated by applicants from an origin across the offices represented in the
bulk file. That volume combines at least two components: the scale of
underlying patenting and the breadth of the office network through which
protection is sought. For questions about documentary reach, filing routes or
the international workload associated with an origin, the combination may be
the object of interest. For comparisons of inventive output, it is a
confounder.

The empirical range shows why this is more than a terminological preference.
A common multiplier would preserve origin ratios even while making every
level too large. The measured proxy is not common. The median pair differs
by 1.57; China and Switzerland sit at 1.05 and 4.48, a 4.3-fold difference
at the sample ends. Office summing therefore embeds an origin-specific filing
profile. Much of that profile is arithmetic: where the home office is the
cell-wise leader, $D=1/H$. The remainder is structure. Observed $D$ lies
below $1/H$ because foreign filings concentrate in a few large offices.
Partial Spearman of $+0.413$ after controlling for $H$ is the evidence that
the high tail remains an EPC pattern rather than a generic
international-orientation gradient; the matched-$H$ interleave (New Zealand
beside Sweden) shows that EP-as-leader is not required for that tail. The
pooled median of 2.27 is a mixture: 35 of the 68 are EPC contracting-state
origins, and the group medians on the measurable-home 58 are the more stable
summary. The
direction of cell coarsening is arithmetic---coarser cells cannot lower
$D$---and the magnitude is the finding: origin-only $D$ reaches 1.05 / 2.54 /
5.77, which is the office-sum divided by the largest single named office.

This does not imply that the cell-wise leading office supplies the true
invention count. China's factor of 1.05 says that other named offices add
little relative to the leading-office baseline in these cells. It does not
say that the Chinese home-office total contains one record per invention.
The baseline can include multiple procedural publications, whole-counted
technology assignments and other features of the source. Likewise,
Switzerland's 4.48 is not an exact average family size. The diagnostic
compares concentration under one consistent aggregation rule; it does not
recover hidden record linkage.

The PCT result adds a separate source of non-comparability. Office
\texttt{**} tracks published PCT applications by origin, but its share is not
constant across origins. In the worked pair, Singapore's 20.2\% is far above
China's 3.9\%. A total formed by adding that route to all named offices would
therefore alter the comparison again. The correct treatment depends on the
question, but the choice cannot remain implicit.

The family analysis then establishes a second distinction. Moving from
office publications to patent families addresses the multi-office assignment
problem because related documents are grouped under a declared family rule.
It does not decide whether domestic-only families belong in the construct.
The China--Singapore ratio in 2022 is about five hundred for all families and
about twenty for the foreign-oriented subset. Neither ratio invalidates the
other. They describe different geographic scopes, and their distance reflects
the contrast between 3.6\% and 91.9\% foreign extension in the two portfolios.

\subsection{Using the file}\label{sec:using}

For a patent-information practitioner, the safest workflow begins before
aggregation, with a sentence that names the intended construct. If the
question concerns publications handled by a specified office, the office
dimension should be fixed and the resulting statistic labelled as
publications. If the question concerns an origin's use of the international
system, a multi-office sum or the PCT layer may be informative, provided the
unit remains publication assignments and the routes are reported separately.
If the question concerns inventions, the analysis needs a family product or
document-level data from which an explicit family definition can be built.

This sequence is especially important in reusable dashboards and country
rankings. A code path that groups by applicant origin and sums every remaining
numeric column can be technically correct yet conceptually wrong. The error
will not be visible as a missing value or an implausible total. Fig.~\ref{fig:ranks}
shows the consequence: equally defensible rules reorder the same origins.
Metadata and method notes therefore need to travel with the exported table,
not remain in the analyst's private understanding.

When only this file is available, the cell-wise leading office is closer to
families than the office-sum over the seven-economy overlap, all non-EPC
origins filing through national rather than regional-first routes. The
validation excludes the EPC structural case. Singapore is the only low-$H$
analogue in that overlap and has the largest genuine error. Four EPC family
exports were not available. The recommendation is therefore scoped to this
overlap; it is not stated for regional-first filing systems. $D$ predicts the
office-sum's distance from the leading office by construction. There is no
evidence at $n=7$ that it predicts the error a user would incur, so it is not
a validated warning light for that proxy error. It does not support dividing
an office-sum by $D_i$ and calling the result families.

Users should also preserve the distinction between the two \texttt{**}
tokens in processing code and documentation. Dropping both with one
unexplained filter would conceal why observations were removed. Origin
\texttt{**} cannot support country attribution. Office \texttt{**}, under
the interpretation used here, is analytically meaningful as the PCT layer.
A reproducible transformation should name the column and reason for each
treatment.

Finally, a family total should carry its scope in the table heading. ``Patent
families by applicant origin'' is incomplete if the reader cannot tell whether
domestic-only families are included. Indicator 6a is suitable when the full
origin portfolio is intended. Indicator 7 is suitable when the construct is
internationally extended inventions. International families, including
triadic constructions, can improve comparability for some research questions
but also select on filing strategy and market reach
\citep{dernis2004triadic,dechezlepretre2017families}. The scope should follow
the question rather than being treated as a generic quality filter.

\subsection{Limitations}\label{sec:limitations}

The first limitation concerns identification of the PCT layer. Origin-level
office-\texttt{**} volumes track published PCT applications by origin
(Spearman $+0.991$ over the twenty largest named origins in 2020), which is
the test an unknown-office bucket would fail. A file-specific dictionary that
explicitly maps this token has still not been located. All PCT-layer findings
are therefore stronger than a $Q$--$H$ correlation but still without a
dictionary confirmation. A different documented meaning would require the
layer to be renamed and reinterpreted.

Second, the duplication factor is intentionally modest in what it claims.
The aggregate bulk file lacks identifiers linking related publications across
offices, so $D_i$ cannot distinguish an international family from unrelated
documents that happen to share a cell. It cannot identify continuations,
procedural publications or family structures. Its denominator is a cell-wise
leading-office baseline, not a proven count of inventions, and 4.3 is a
fold-difference between the proxy's sample ends rather than a deduplication
bound. $D$ also confounds office dispersion with the temporal spread of
publication; the origin--field grain bounds that timing component.

Third, the comparison with family indicators contains seven economies, all
non-EPC origins filing through national rather than regional-first routes.
The leading-office comparison discloses a scoped practical recommendation,
but the exercise is not a validation study across the full sample and it
excludes the EPC structural case the paper itself identifies. Family exports
for four EPC origins were not available. In particular, the rank check does
not establish a conversion from the proxy to family counts, and $D$ is not a
validated warning light for the error a user would incur. More overlapping
origins or record-level family links would be needed to test general
calibration.

Fourth, whole counting presents a documentation tension. The bulk-file values
are integers and the Schmoch concordance supports assignment to multiple
technology fields, while WIPO's general Data Center description refers to
fractional counting for technology indicators \citep{schmoch2008,wipoabout}.
This article reports the file as observed and retains a consistent
field-cell construction. It cannot resolve whether the general documentation
was intended to describe this exact download. This uncertainty reinforces the
need to call the 50.2 million total assignments rather than documents.
Whole-versus-fractional counting and the identification of office \texttt{**}
are both questions about how one PATSTAT compilation was run
\citep{wipoabout}.

Fifth, WIPO attributes origin to the residence of the first-named applicant
and ignores co-applicants \citep{wipoabout}. Multi-jurisdiction filings by
multinational groups can therefore carry different first-named applicants at
different offices, so the same invention may be attributed to different
origins in different office cells. That attribution rule places
micro-state and offshore-domicile origins at the extremes of the structural
measures this file produces. Antigua and Barbuda has $D=4.10$, origin-only
$D=5.77$, 1,592 named-office publications and Japan as origin-level leading
office; its home office is unmeasurable, and it is the all-68 non-EPC
maximum and the largest rank movement among the 68 (50th versus 65th).
Monaco has the largest labelled-versus-computed $D$ gap, 2.294, with $D=2.56$,
the EPO as origin-level leader, a measurable home office and EPC membership.
Iran has the sample-maximum PCT-layer share, 54.4\%, and an unmeasurable home
office; it belongs in this class rather than in the PCT-layer result as a
confirmed filing behaviour (Section~\ref{sec:pct-result}).
None of these is an inventing-economy profile in the ordinary sense; they
follow from first-named-applicant attribution and from unmeasurable home
offices. The 500-publication threshold
does not screen this, because it screens volume rather than attribution
meaningfulness. Singapore remains a hub case whose publications are
distributed across foreign offices, but it is no longer the obvious
illustration of the attribution rule.

Sixth, the family tail is subject to right truncation. Families are dated by
earliest filing year, but foreign members can appear later. The 2022
foreign-oriented shares and ratio are therefore provisional, and the 2018
anchor is reported as a more mature comparison. This issue does not affect
the 2018--2022 publication-window range in the same way because those
observations are dated by publication year, but it limits direct annual
alignment between products.

Finally, the home office is unmeasurable for 10 of the 68 economies. Retaining
their observed-office duplication factors preserves sample coverage, but a
missing home-office block may omit the true cell-wise leader and raise the
proxy. Saudi Arabia's 32 origin-resolved home-office publications place it at
that boundary rather than clearly inside the measurable set
(Section~\ref{sec:dup-result}). The separate median for the 58 measurable
economies and the exclusion of missing shares from Fig.~\ref{fig:dup} make
that limitation visible. Neither the threshold sensitivity, the cell-grain
check nor the reported medians can replace origin-resolved source coverage.

\section{Conclusions}\label{sec:conclusions}

Three findings hold without the scoped qualifications that follow. The
office-sum is not an invention count: relative to the cell-wise leading
office, $D$ ranges from 1.05 to 4.48 across 68 economies (median pairwise
ratio 1.57), and it is not a common multiplier. Where the home office leads,
that range is the identity $D=1/H$; where it does not, $D$ stays far below
$1/H$ because foreign filings concentrate in a few large offices. Origin
rankings are not invariant to defensible counting rules, demonstrated on the
top 25.

The WIPO patent-publications-by-technology bulk file supports useful
cross-country work when that assignment structure is respected. Partial
Spearman of $D$ with EPC status controlling for $H$ is $+0.413$, so the high
tail remains an EPC pattern after home-office share is held; EP-as-leader is
not required for it. Office \texttt{**} tracks published PCT applications by
origin (Spearman $+0.991$) and should be reported according to the question
being asked. Over the seven-economy overlap---all non-EPC origins filing
through national rather than regional-first routes---the cell-wise leading
office is closer to families than the office-sum (mean absolute relative
error 0.101, or 0.069 excluding the year-anchor case, versus 0.803). That
recommendation is not validated for the EPC structural case; Singapore is the
only low-$H$ analogue in the overlap and has the largest remaining error
(0.22). Once families are used, geographic scope must still be defined: all
families and foreign-oriented families produce different but legitimate
China--Singapore comparisons (about 500 versus about 20). Interpretable
patent statistics therefore require three explicit choices---counting unit,
treatment of the PCT layer and family scope---before the country totals are
compared.

\section*{Data availability}

The bulk file and the family export are public. The derived tables, figure
files and machine-readable record of numeric claims are archived with the
analysis pipeline at
\url{https://github.com/samiurk70/WIPO-Patent-Research}.
Per-economy home-office share, PCT-layer share and duplication factor for the
68 economies are in \texttt{counting\_generality.csv} in that archive.
A Zenodo snapshot will be linked here when minted.

\section*{Declaration of competing interest}

The authors declare that they have no known competing financial interests or
personal relationships that could have appeared to influence the work
reported in this paper.

\section*{Declaration of generative AI and AI-assisted technologies in the
manuscript preparation process}

During the preparation of this work the authors used an AI-assisted editor in
order to organise the draft from the authors' analysis pipeline and to match
the Elsevier CAS template. After using this tool, the authors reviewed and
edited the content as needed and take full responsibility for the content of
the published article.

\section*{Acknowledgements}

The authors thank the maintainers of the WIPO IP Statistics Data Center. The
patent-family series was exported by hand from that interface; the remaining
inputs were retrieved by script.

\printcredits

\bibliographystyle{elsarticle-num}
\bibliography{refs}

\end{document}
